%%=============================================================================
%% Conclusie
%%=============================================================================

\chapter{Conclusie}%
\label{ch:conclusie}

% TODO: Trek een duidelijke conclusie, in de vorm van een antwoord op de
% onderzoeksvra(a)g(en). Wat was jouw bijdrage aan het onderzoeksdomein en
% hoe biedt dit meerwaarde aan het vakgebied/doelgroep? 
% Reflecteer kritisch over het resultaat. In Engelse teksten wordt deze sectie
% ``Discussion'' genoemd. Had je deze uitkomst verwacht? Zijn er zaken die nog
% niet duidelijk zijn?
% Heeft het onderzoek geleid tot nieuwe vragen die uitnodigen tot verder 
%onderzoek?

Dit hoofdstuk beantwoordt de centrale onderzoeksvraag en de deelvragen uit
Hoofdstuk~\ref{ch:inleiding}, op basis van de theoretische bevindingen uit
Hoofdstuk~\ref{ch:stand-van-zaken} en de empirische resultaten uit
Hoofdstuk~\ref{ch:resultaten-en-evaluatie}. Vervolgens wordt de bijdrage van dit
onderzoek aan het vakgebied toegelicht, gevolgd door een kritische reflectie op de
gehanteerde aanpak en een aanzet voor toekomstig onderzoek.
 
\section{Antwoord op de centrale onderzoeksvraag}%
\label{sec:conclusie-hoofdvraag}
 
De centrale onderzoeksvraag luidde: \textit{hoe kan een verbinding worden opgezet met
een eigen fileserver via een radioverbinding, en wat zijn de technische randvoorwaarden
en maximale afstandsbeperkingen voor een werkende implementatie?}
 
Het antwoord is genuanceerd, en net die nuance vormt de kern van de bijdrage van dit
onderzoek. Een architectuur die een private fileserver via SFTP over een
point-to-point-straalverbinding ontsluit, is technisch haalbaar en presteert binnen het
optimale bereik (0--15~km) ruimschoots binnen de vooropgestelde eisen: het gemiddelde
packet loss onder normale kanaalomstandigheden bedroeg er slechts 0{,}35\%, ruim onder
de norm van 2\% (REQ-03, Sectie~\ref{sec:res-interpretatie}). Vanaf de horizonlimiet
(15--30~km) schiet diezelfde kanaalkwaliteit echter reeds structureel tekort (gemiddeld
5{,}27\% verlies), en voorbij het bereik dat in de requirementsanalyse als onbetrouwbaar
werd bestempeld (30+~km) is die tekortkoming vrijwel universeel (20{,}38\%). De
"maximale afstandsbeperking" uit de onderzoeksvraag is dus geen enkel harde grenswaarde,
maar een geleidelijke, kwantitatief in kaart gebrachte overgang: de architectuur is
\emph{geschikt} binnen ongeveer 15~km, \emph{grensgeval} tot ongeveer 30~km, en
\emph{ongeschikt} daarbuiten voor permanente, betrouwbare bestandstoegang.
 
Eén element van de architectuur bleek daarbij opmerkelijk consistent, ongeacht de
afstand: het vermogen om een bestandsoverdracht te hervatten na een verplichte,
wettelijk voorgeschreven kanaalonderbreking door Dynamic Frequency Selection (FR-06,
Sectie~\ref{subsec:dfs}). Op alle drie de afstandscategorieën herstelde de combinatie
van SFTP en TCP zich betrouwbaar van zo'n onderbreking, wat aantoont dat deze
protocolkeuze een robuuste basis vormt, onafhankelijk van de onderliggende
kanaalkwaliteit.
 
\section{Antwoord op de deelvragen}%
\label{sec:conclusie-deelvragen}
 
\paragraph{Welke radioprotocollen zijn geschikt voor bestandsoverdracht?} De
literatuurstudie (Hoofdstuk~\ref{ch:stand-van-zaken}) toonde aan dat klassieke
amateurradioprotocollen zoals AX.25 en Winlink, hoewel robuust en breed inzetbaar voor
noodcommunicatie, primair ontworpen zijn voor tekstgebaseerde berichten met een lage
bandbreedte, en zich niet lenen tot het efficiënt overdragen van bestanden van
enige omvang. Native IP-over-radio -- het tunnelen van gewoon IP-verkeer over een
radioverbinding, waardoor bestaande, standaard bestandsoverdrachtsprotocollen hergebruikt
kunnen worden -- bleek de geschiktere aanpak. Binnen die protocollen is specifiek SFTP
gekozen boven bijvoorbeeld SMB, omwille van zijn beperkte "chattiness" (minder
gevoelig voor latency- en jitterschommelingen, Sectie~\ref{sec:req-tussentijdse-conclusie})
en zijn ingebouwde authenticatie en versleuteling.
 
\paragraph{Welke multiplextechnieken zijn van toepassing, en welke beperkingen leggen
zij op?} OFDM werd in de literatuurstudie geïdentificeerd als de meest geschikte
multiplextechniek voor instabiele radiokanalen, dankzij zijn robuustheid tegen
multipath-fading en looptijdverschillen (Sectie~\ref{sec:Multiplextechnieken in radiogebaseerde datacommunicatie}).
Deze keuze kon in de proof-of-concept echter niet rechtstreeks empirisch geverifieerd
worden: de gebruikte netwerkemulatie (Hoofdstuk~\ref{ch:proof-of-concept}) simuleert het
gedrag van de fysieke laag op basis van statistische parameters (vertraging,
doorvoersnelheid, burstverlies), maar bootst geen daadwerkelijke modulatie of
multiplexing na. Dit blijft dan ook een aanbeveling voor verder onderzoek
(Sectie~\ref{sec:conclusie-verder-onderzoek}).
 
\paragraph{Wat is de maximale praktische afstand?} Dit is de deelvraag waarop dit
onderzoek het meest concrete antwoord biedt, samengevat in
Sectie~\ref{sec:conclusie-hoofdvraag} hierboven en in detail onderbouwd in
Tabel~\ref{tab:resultaten-per-afstand} en Tabel~\ref{tab:res-toetsing}. In tegenstelling
tot een eenmalige meting op een enkele afstand, geeft de statistische onderbouwing over
300 trials (100 per afstandscategorie, met telkens een willekeurige weers- en
DFS-realisatie) een betrouwbaar beeld van de \emph{spreiding} van de haalbare prestaties,
niet enkel van een toevallig gunstig of ongunstig moment.
 
\paragraph{Welke minimale infrastructuur en configuratie zijn vereist?} Hoofdstuk~\ref{ch:proof-of-concept}
beschrijft de minimale opstelling in detail: een SFTP-server met een chroot-omgeving
voor de gebruiker, een netwerkbrug die de radioverbinding representeert, en een client
met standaard SFTP-, meet- en diagnosetools. De volledige opstelling -- van installatie
tot statistisch onderbouwde meting -- is bovendien volledig geautomatiseerd, wat het
antwoord op deze deelvraag ook onmiddellijk reproduceerbaar en overdraagbaar maakt naar
andere onderzoekers of praktijksituaties.
 
\section{Bijdrage aan het vakgebied en meerwaarde voor de doelgroep}%
\label{sec:conclusie-bijdrage}
 
De literatuurstudie in Hoofdstuk~\ref{ch:stand-van-zaken} identificeerde een lacune in
de bestaande vakliteratuur: onderzoek naar radiogebaseerde communicatie richt zich
overwegend op noodcommunicatie, professionele draadloze backhaul, of berichtenverkeer,
maar niet specifiek op de toegankelijkheid van een persoonlijke of kleinschalige
fileserver. Dit onderzoek vult die lacune op met een kwantitatief, per afstand
gedifferentieerd antwoord, dat rechtstreeks bruikbaar is voor de in
Sectie~\ref{sec:probleemstelling} geïdentificeerde doelgroep: maritieme operatoren,
hulpverleningsorganisaties, onderzoekers in afgelegen gebieden en radioamateurs kunnen
uit dit onderzoek een concreet, onderbouwd beeld meenemen van tot welke afstand een
dergelijke opstelling betrouwbaar inzetbaar is, en welke architecturale keuzes (SFTP,
DFS-conforme apparatuur) daarbij aangewezen zijn.
 
Daarnaast levert dit onderzoek een methodologische bijdrage die verder reikt dan de
specifieke onderzoeksvraag: de volledig geautomatiseerde, statistisch onderbouwde
testmethodologie uit Hoofdstuk~\ref{ch:proof-of-concept} -- met een Gilbert-Elliot-
burstverliesmodel voor weersvariatie en een stochastische DFS-simulatie, herhaald over
honderden onafhankelijke trials -- is generiek genoeg om ook voor andere
radiogebaseerde of draadloze netwerkonderzoeken te worden hergebruikt. Waar veel
kleinschalige proof-of-concept-onderzoeken zich beperken tot een handvol
puntmetingen, toont dit onderzoek aan dat een statistisch representatieve steekproef,
ook binnen de tijdspanne van een bachelorproef, haalbaar en waardevol is.
 
\section{Kritische reflectie}%
\label{sec:conclusie-reflectie}
 
Het algemene patroon van afnemende kanaalkwaliteit met toenemende afstand was verwacht:
zowel de wet van Shannon-Hartley (Hoofdstuk~\ref{ch:stand-van-zaken}) als de
kwalitatieve inschatting in de requirementsanalyse (Sectie~\ref{sec:req-tussentijdse-conclusie})
voorspelden dat de 30+~km-categorie buiten het betrouwbare toepassingsgebied zou vallen.
De mate waarin dit vermoeden nu kwantitatief bevestigd wordt -- 99\% van de trials
overschrijdt er de norm, tegenover slechts 15\% bij 0--15~km -- was echter scherper dan
op voorhand kon worden aangenomen. Minder verwacht was het onderscheid dat de
DFS-opsplitsing binnen de 15--30~km-categorie blootlegde: daar bleek de
REQ-03-overschrijding \emph{niet} voornamelijk het gevolg van de verplichte
DFS-onderbrekingen (zoals bij 0--15~km), maar van de courante kanaalkwaliteit zelf. Dat
onderscheid was pas zichtbaar dankzij de expliciete opsplitsing van de resultaten, en
onderstreept het belang van een voldoende gedetailleerde statistische analyse boven een
enkel, samengevat gemiddelde.
 
De belangrijkste beperking van dit onderzoek is methodologisch van aard, en werd reeds
uitgebreid toegelicht in Sectie~\ref{subsec:meth-bijstelling-fase3}: de voorziene
fysieke 5.6~GHz airMAX-straalverbinding kon niet worden aangekocht, waardoor de volledige
proof-of-concept via netwerkemulatie gerealiseerd werd. De gerapporteerde resultaten
tonen daardoor aan of de gekozen architectuur functioneel en performant is onder
\emph{verwachte}, uit officiële specificaties afgeleide kanaalcondities, maar
bevestigen niet de effectieve RF-prestaties -- antennealignment, werkelijke
signaalpropagatie, concrete weersinvloed -- van een fysieke link in het veld. De
gebruikte Gilbert-Elliot-parameters en DFS-kans zijn plausibele, uit de literatuur en
productspecificaties afgeleide schattingen, geen gemeten grootheden. Deze beperking
raakt met name de deelvraag over multiplextechnieken (Sectie~\ref{sec:conclusie-deelvragen}),
die in deze opstelling niet op fysiek niveau getoetst kon worden.
 
Een tweede, kleinere kanttekening betreft de asymmetrie in meetintensiteit tussen de
requirements: REQ-03 (packet loss) en FR-06 (DFS-herstel) werden dankzij hun
niet-deterministische aard over 300 trials statistisch onderbouwd, terwijl REQ-01
(doorvoersnelheid) en REQ-02 (latency/jitter) via een beperkter aantal metingen per
afstandscategorie werden gevalideerd (Sectie~\ref{sec:sim-uitvoering}). Een even
uitgebreide statistische onderbouwing van die laatste twee eisen zou het totaalbeeld
verder versterken.
 
\section{Aanbevelingen voor toekomstig onderzoek}%
\label{sec:conclusie-verder-onderzoek}
 
Op basis van de bevindingen en beperkingen hierboven worden de volgende richtingen voor
vervolgonderzoek voorgesteld:
 
\begin{itemize}
  \item \textbf{Fysieke veldvalidatie.} Zodra geschikte 5.6~GHz airMAX-apparatuur (of een
        vergelijkbaar alternatief) beschikbaar is, zou een fysieke herhaling van de
        drie afstandscategorieën toelaten om de in dit onderzoek gebruikte
        Gilbert-Elliot- en DFS-parameters te kalibreren aan werkelijke veldmetingen, en
        om architecturale aspecten (antennealignment, effectieve modulatie) te
        valideren die in een netwerkemulatie noodgedwongen buiten beschouwing bleven.
  \item \textbf{Uitbreiding van de statistische onderbouwing naar REQ-01 en REQ-02.} Een
        vergelijkbare batchmethodologie als voor packet loss (Hoofdstuk~\ref{ch:proof-of-concept})
        zou ook voor doorvoersnelheid en latency toegepast kunnen worden, voor een
        even robuust onderbouwd beeld van alle vier de netwerktechnische requirements.
  \item \textbf{Alternatieve frequentiebanden en vermogensklassen.} Dit onderzoek
        beperkte zich, om regelgevende redenen, tot de licentievrije 5.6~GHz-limieten
        (1~W EIRP). Een vergelijkende studie met het gelicentieerde amateursegment
        (tot 200~W PEP, Sectie~\ref{sec:req-tussentijdse-conclusie}) of met de 6m-band
        uit Hoofdstuk~\ref{ch:stand-van-zaken} zou kunnen uitwijzen of een groter bereik
        haalbaar is binnen een andere regelgevende context.
  \item \textbf{Multi-hop-oplossingen voor het onbetrouwbare bereik.} Voor de
        30+~km-categorie, waar dit onderzoek aantoonde dat een directe verbinding niet
        volstaat, zou onderzocht kunnen worden of een keten van meerdere, kortere
        straalverbindingen (relais) een uitkomst biedt -- elke afzonderlijke schakel
        zou dan binnen het in dit onderzoek als betrouwbaar aangetoonde bereik van
        0--15~km kunnen blijven.
  \item \textbf{Reële, gemeten weersdata.} In plaats van willekeurig getrokken
        Gilbert-Elliot-parameters binnen een plausibele bandbreedte, zou toekomstig
        onderzoek de simulatie kunnen voeden met effectief gemeten meteorologische
        gegevens (regenintensiteit, luchtvochtigheid) gekoppeld aan gelijktijdige
        signaalmetingen, voor een nog realistischer verband tussen weer en
        kanaalkwaliteit.
\end{itemize}
 
Samengevat toont dit onderzoek aan dat radiogebaseerde toegang tot een private
fileserver, binnen een zorgvuldig afgebakend afstandsbereik, een technisch haalbaar en
statistisch onderbouwd alternatief vormt voor de doelgroep uit
Sectie~\ref{sec:probleemstelling} -- met een duidelijk, kwantitatief in kaart gebracht
omslagpunt waarop toekomstig onderzoek verder kan bouwen.

